\documentclass[aps,preprint,showpacs,amsmath,amssymb,prc,floatfix,nofootinbib]{revtex4-1}
\usepackage{color}
\usepackage{natbib}
\usepackage{times}
\usepackage{inputenc}
\usepackage{graphicx}
\usepackage[colorlinks,citecolor=blue,linkcolor=red,anchorcolor=blue,filecolor=blue,urlcolor=blue]{hyperref}
\evensidemargin 0.0cm
\topmargin -1.7 cm
\oddsidemargin -1.0 cm
\textheight  21.0 cm
\textwidth 18.5cm
% \usepackage{fancyhdr}
% \pagestyle{fancy}
% \fancyhf{}
% \fancyfoot[LE,RO]{\thepage}

%%%%%%%%%%%%%%%%
\begin{document}
%%%%%%%%%%%%%%%%
\begin{center}
    \textbf{\underline{Referee Comments}}
\end{center}
%%%%%%%%%%%%%%%
This thesis titled ``Impacts of dark matter interaction on nuclear and neutron
star matter within the relativistic mean-field model" and authored by Harish
Chandra Das, is about the effect of dark matter inside the neutron star (NS)
and to explore its effects on various NS properties. This appears to be quite a
unique topic, particularly when the equation of state (EoS) of NS matter and
underlying interactions is still quite uncertain. Indeed the thesis is based on the seven (7) papers published mostly in internationally well reputed journals. In
all the papers, the candidate is the first/lead author. The thesis merits a Ph.D.
degree.
\\
\\
The thesis is composed of seven chapters including Introduction (Chapter 1) and Summary and Conclusion (Chapter 7). Chapter 2 revises in detail the relativistic mean-field models and their applications, which is the basic building block of the underlying EoS. Chapter 3 recalls some basics of dark matter interaction and its effects on nuclear matter properties. Therefore, Chapters 2 and 3 serve as the basic ingredients of the main exploration of the thesis. Chapter 4 initiates understanding the impacts of dark matter on isolated and static/rotating NS properties. The candidate first constructs the EoSs and then obtains the underlying mass-radius relations. He shows that while the maximum mass of a NS could be as high as $2.7M$, it decreases with the inclusion of dark matter. He also explores other general relativistic properties. Chapters 5 and 6 explores dark matter effects on properties of the oscillating NSs/hyperon stars, and binary NS properties. The latter is more relevant to gravitational wave astronomy. The author has calculated the dark matter admixed binary NS properties with RMF and E-RMF EoSs. The shape of the binary NS is deformed due to the mutual interactions between the components. The rate of deformations can be measured by various Love numbers and their corresponding deformabilities.
\\
\\
However, I have the following comments and suggestions for amendments. Overall, however, the thesis deserves a Ph.D., and I recommend for the same, particularly once the all corrections are amended in the thesis.
\\
\\
{\bf Comment-1:} Chapter 1.1: The mentioned magnetic field for NS is a couple of orders of magnitude lower than the current understanding, if it is surface field. But internally NS field could be even $10^{18}$ G.
\\
{\bf Response:} \textcolor{blue}{Thanks for the suggestions. The magnetic field at the surface of the NS is $\sim 10^{15}$ G, while it is roughly $10^{18}$ G at the core. It is approximately $10^8-10^{15}$ times stronger than the Earth's magnetic field.}
\\
\\
{\bf Comment-2:} Chapter 1.1.2: I think the candidate mixes up between Chandrasekhar limit of white dwarf with neutron drip. At a higher density inverse-beta decay takes place.
\\
{\bf Response:} \textcolor{blue}{Thanks for the comments. We removed the word ``Chandrasekhar limit" for this confusion.}
\\
\\
{\bf Comment-3:} Many grammatically incomplete/incorrect sentences are there: ``The millisecond pulsar (MSP) J0740+6620, having mass $M = 2.14^{+0.10}
_{-0.09} M_\odot$ with 68.3\% credibility interval (CI) [17] using the same Shapiro delay measurement.", ``The simultaneous measurement of mass and radius of the PSR J0030+0451 by the NASA Interior Composition Explorer (NICER), a telescope fixed in the international space station.", ``...GW170817 event opens the possibility that the NS can be accumulated some amount of DM.", ``Hence, the Lagrangian density of the system is the addition of NS and DM (see Chapter 3.2)."
\\
{\bf Response:} \textcolor{blue}{We are sorry for this. These lines are corrected in the revised version.}
\\
\\
{\bf Comment-4:} Either you should write NS (once NS is defined) or `neutron star' throughout. So many notations are not defined, e.g., NM (I assume nuclear matter). 
\\
{\bf Response:} \textcolor{blue}{Thanks for this suggestion. We changed the `neutron star' as NS in the entire revised manuscript. We have defined other also.}
\\
\\
There are many statements here and there which either seem to be incorrect or require better way of rewriting or require explanation. Some examples are given below:
\\
\\
{\bf Comment-5:} Below equation (2.17), $E_j^*$ is said to be effective energy density. but dimensionally based on its definition it is energy, not energy density, unless the candidate means something else in the definition.
\\
{\bf Response:} \textcolor{blue}{Thanks for this suggestions. We have modified in the revised version.}
\\
\\
{\bf Comment-6:} Below equation (2.21), it is mentioned that ``the third component $<T_{ii}>$ determines the pressure of the system". Why ``third component", not spatial component, and if it is indeed, why not $<T_{33}>$? 
\\
{\bf Response:} \textcolor{blue}{Sorry for this misunderstanding. We have taken care in the revised version.}
\\
\\
{\bf Comment-7:} In the beginning of Chapter 2.2.3, it is mentioned that ``The neutron decays to proton, electron, and anti-neutrino inside the NS". However, we know that for density above $10^7$ gm/cc inverse-beta decay happens and the neutron
star density is several orders of magnitude higher (hence neutron drip happens,
whose threshold is $3\times10^{11}$ gm/cc). Hence, the inverse-beta-decay is expected, unless the author means it for very outer region. This should be then explicitly mentioned.
\\
{\bf Response:} \textcolor{blue}{We are thankful to referee for this question. The inverse $\beta$--decay process occurred to maintain the charge neutrality condition, mainly at the core part where neutron stars drifted from the nuclei. We have modified this in the revised version.}
\\
\\
{\bf Comment-8:} The statement ``The Universe is made of approximately 85\% of dark
matter." is highly confusing. I understand what the candidate means: 85\% of
total matter. However, the energy is also a form of matter and universe is mostly
consisting of dark energy (about 68\%). In that sense, only 27\% is dark matter.
Either the author should mention it explicitly or change the percentage.
\\
{\bf Response:} \textcolor{blue}{It is corrected in the revised version. Around 27\% of the matter in the universe is believed to be composed of an unobservable and speculative type of matter known as dark matter.}
\\
\\
{\bf Comment-9:} Above Fig. 2.12, the description of the figure is in a way as if only IOPB-I and FSUGarnet EoSs are there. However, NL3 and BigApple saturate
at a much higher value (in fact former does not seem to saturate). Hence, this
part has to be modified. Also, why they don’t obey c/3 limit, no comment?
\\
{\bf Response:} \textcolor{blue}{Thanks for this question. NL3 and BigApple EOSs are stiffer in nature compared to others. Therefore, the derivative (or slope) of pressure with respect to energy density is larger than other EOSs such as IOPB-I and FSUGarnet. The $c/3$ limit is only valid at the higher density part where quarks play a role. In this case, we only consider the nucleons. Therefore, QCD conformal limit doesn't obey the considered parameter sets as shown in Fig. \ref{fig:cs2_NS_1}.}
%%%%%%%%%%%%%%
\begin{figure}
\centering
\includegraphics[width=0.6\columnwidth]{Figures/Chapter-2/CS.pdf}
\caption{The speed of the sound as the function of baryon density for different RMF models. The dashed magenta lines represent the causality bound ($C_s^2=1$) and QCD conformal limit ($C_s^2=1/3$).}
\label{fig:cs2_NS_1}
\end{figure}
%%%%%%%%%%%%%%
\\
\\
{\bf Comment-10:} In Chapter 3.3.2 (paragraph above 3.3.3), initially the conjecture of $C_s \leq c/\sqrt{3}$ was argued, which however does not preserve in Fig. 3.5, of what no further explanation, except that causality, preserved. This should be expanded.
\\
{\bf Response:} \textcolor{blue}{Thanks for this question. The RMF formalism is always causal in nature. However, the density at the core of the neutron star is so high that the nuclei break into quarks. In that medium, the sound speed is basically slower in comparison to the nucleonic medium. This is the main tension that the sound speed violates the QCD conformal limit for only nucleonic cases as discussed in {\bf Refs. [PRL 114, 031103 (2015), PRL 122, 122701 (2019)]}. In this study, we include only nucleons. Therefore, it violates that limit in the whole region of the stars. This part is added in the revised version of the thesis.}
%%%%%%%%%%%%%%
\begin{figure}
\centering
\includegraphics[width=0.6\columnwidth]{Figures/Chapter-3/vel.eps}
\caption{The speed of the sound as a function of total baryon density with $k_f^{\rm DM}$ = 0.06 GeV.}
\label{fig:cs2_NS_2}
\end{figure}
%%%%%%%%%%%%%%
\\
\\
{\bf Comment-11:} In Chapter 4, RNS and particularly SNS should be explicitly defined; even if RNS is known to be an Einstein equation solver code, SNS does not seem to be a common word. Also in Fig. 4.3, radius does not mean much when the star is (fast) rotating and polar, and equatorial radii are different. The same goes with other figures reporting the radius of rotating NSs.
\\
{\bf Response:} \textcolor{blue}{Thanks for the suggestions. It is defined in the revised version.}
\\
\\
{\bf Comment-12:} Overall, in the same Chapter as above, a NS mass is shown to be 2.78 solar mass for the NL3 EoS with DM effect. I am a bit (happily) surprised about it when it seems to overcome the issue of hyperon puzzle. Unless, I miss a relevant discussion in the thesis, this point has to be elaborated a bit. Is it the hypothetical pure nuclear matter case?
\\
{\bf Response:} \textcolor{blue}{Thanks for this question. However, we have already discussed the hyperon puzzle in chapter-5 in the introduction part.}
\\
\\
{\bf Comment-13:} There are several similar apparent lose statements in the remaining chapters, what the candidate may like to revise.
\\
{\bf Response:} \textcolor{blue}{We have gone through the each chapters, and modified accordingly with best of our knowledge.}
\\
\\
We have given almost all answers to the examiner's questions and also inserted them in the revised thesis. Now, we think that the thesis is much better, all thanks to the referee, and hopefully, the referee will find the thesis is suitable for the award of Ph. D. degree.
\\
\\
Best Regards
\\
Sincerely yours
\\
H. C. Das
\\
Email: harish.d@iopb.res.in
%%%%%%%%%%%%%%%%%%%%%%%%%
%%%%%%%%%%%%%%%
\end{document}
%%%%%%%%%%%%%%%%%%%%%%%%%%%%%%%%%%% END %%%%%%%%%%%%%%%%%%%%%%%%%%%%%%%%%%%%%%%%%%